\documentclass[11pt,a4paper]{article}

% --- Pakete (Overleaf: pdfLaTeX) ---
\usepackage[utf8]{inputenc}
\usepackage[T1]{fontenc}
\usepackage[ngerman]{babel}
\usepackage{lmodern}
\usepackage{booktabs}
\usepackage{graphicx}
\usepackage{xcolor}
\usepackage{amsmath}
\usepackage[margin=2.4cm]{geometry}
\usepackage[colorlinks=true,linkcolor=black,urlcolor=blue]{hyperref}

% Methodenfarben (wie im Viewer)
\definecolor{fdgreen}{HTML}{2E9E4A}
\definecolor{cncyan}{HTML}{1C9BAB}
\definecolor{unred}{HTML}{D53030}
\definecolor{samviolet}{HTML}{B03CD0}

\title{\textbf{Semantische Segmentierung der Laserlinie}\\[2pt]
\large SAM als Foundation-Modell-Baseline und die Wahl eines spezialisierten U-Nets}
\author{Eisdickenmessung im Enteisungs-Windkanal --- Methodenvergleich}
\date{\today}

\begin{document}
\maketitle

\begin{abstract}
\noindent
Zur Vermessung der Eisdicke wird eine auf die Slat-Vorderkante projizierte
Laserlinie detektiert; ihre senkrechte Verschiebung über der Bogenlänge ergibt
die Eisdicke. Dieses Dokument fasst zusammen, warum das Foundation-Modell
\emph{Segment Anything (SAM)} für diese 1-Pixel-Linie \emph{zero-shot} ungeeignet
ist, welche quantitativen Belege das zeigt, welche Konsequenzen daraus folgen und
welche Modellklassen tatsächlich passen würden. Ergebnis: ein kompaktes, von Grund
auf trainiertes U-Net (semantische Segmentierung dünner Kurvenstrukturen) ist der
angemessene Ansatz.
\end{abstract}

\section{Was ist SAM?}
\emph{SAM} (Segment Anything Model, Meta AI, 2023) ist ein \textbf{Foundation-Modell}
für Bildsegmentierung. Eigenschaften:
\begin{itemize}
  \item \textbf{Promptbar:} Eingabe von Punkten, Box oder grober Maske $\rightarrow$
        Ausgabe einer Segmentmaske.
  \item \textbf{Zero-Shot:} kein aufgabenspezifisches Training nötig.
  \item \textbf{Trainingsbasis:} SA-1B mit ca.\ 11 Mio.\ Bildern und 1 Mrd.\ Masken
        (Alltags-/Naturbilder).
  \item \textbf{Architektur:} schwerer Bild-Encoder (Vision Transformer) +
        Prompt-Encoder + leichter Masken-Decoder. \emph{SAM2} (2024) erweitert auf Video.
\end{itemize}
In dieser Arbeit wurde die kompakte, schnelle Variante \emph{MobileSAM} (über
\texttt{ultralytics}) als Baseline eingesetzt.

\section{Warum SAM für die Laserlinie ungeeignet ist}
\label{sec:warum}
Vier Mechanismen greifen ineinander:
\begin{enumerate}
  \item \textbf{Interne Herunterskalierung auf ca.\ $1024\times1024$.} SAM skaliert
        jedes Bild auf eine feste Arbeitsauflösung. Der Crop ($1766\times2033$) wird
        verkleinert, wodurch die nur \mbox{1--2\,px} breite Laserlinie
        \emph{sub-pixel} wird und mit dem Hintergrund verschmilzt --- sie liegt
        \emph{unter} SAMs Arbeitsauflösung.
  \item \textbf{Auf Flächen/Objekte trainiert, nicht auf 1-px-Kurven.} Der
        Masken-Decoder erzeugt zusammenhängende \emph{Flächen}. Bei Prompt-Punkten
        auf der Linie greift SAM die umgebende helle Slat-Fläche.
  \item \textbf{Domänen-Fehlpassung.} SA-1B besteht aus Alltagsfotos; ein
        Graustufen-Laserbild im Windkanal ist weit entfernt davon.
  \item \textbf{Falsche Ausgabe-Geometrie.} Benötigt wird eine \emph{1D-Kurve},
        SAM liefert \emph{2D-Flächen} --- ein grundsätzlicher Mismatch.
\end{enumerate}

\begin{table}[h]
  \centering
  \caption{SAM zero-shot auf Serie 260402-174444 (MobileSAM). Der Laser bedeckt
  real $<1\,\%$ des Crops.}
  \label{tab:sam}
  \begin{tabular}{lll}
    \toprule
    \textbf{Betriebsart} & \textbf{Ergebnis} & \textbf{Bewertung}\\
    \midrule
    Punkt-Prompts (auf der Linie) & Maske deckt $\sim$34\,\% des Crops & Riesen-Blob \\
    Automatisch (``alles'')       & 6--7 Flächen-Regionen            & keine ist die Linie \\
    Enger Crop (beste Chance)     & Maske deckt $\sim$74\,\% des Fensters & füllt Fläche statt Linie \\
    \bottomrule
  \end{tabular}
\end{table}

% Zum Kompilieren die Overlays aus laser_v2/sam/ nach Overleaf hochladen.
\begin{figure}[h]
  \centering
  % \includegraphics[width=0.42\textwidth]{engcrop_0040.png}
  \fbox{\parbox[c][4.6cm][c]{0.7\textwidth}{\centering\itshape
    Abbildung: \texttt{engcrop\_0040.png} hochladen und obige Zeile einkommentieren.\\[4pt]
    Enger Crop: die grünen Prompt-Punkte liegen exakt auf dem Laser,
    die rote SAM-Maske bedeckt jedoch die gesamte helle Fläche ($\sim$74\,\%).}}
  \caption{SAM greift selbst zugezoomt die Fläche statt der dünnen Linie.}
  \label{fig:engcrop}
\end{figure}

\section{Konsequenzen}
\label{sec:konsequenzen}
\begin{enumerate}
  \item \textbf{Ein großes Foundation-Modell ist nicht automatisch gut für eine
        spezialisierte Messaufgabe.} Größe/Allgemeinheit schlägt die
        Aufgabenspezifik bei 1-px-Präzision nicht.
  \item \textbf{Empirische Rechtfertigung des eigenen U-Nets from scratch.} Das
        Versagen von SAM ist der Beleg, dass Spezialisierung die richtige Wahl ist.
  \item \textbf{Rechtfertigung der klassischen Verfahren.} Frame-Differencing und
        Canny arbeiten in voller Auflösung (sub-pixel) und treffen die Linie dort,
        wo SAM (feste 1024\,px) scheitert.
  \item \textbf{Kernlektion: Auflösung ist bei dünnen Strukturen entscheidend.}
        Für 1-px-Objekte muss nativ auflösend gearbeitet werden (Patches/Vollbild),
        nicht herunterskaliert --- genau das leistet das eigene U-Net.
  \item \textbf{SAM ist nicht generell schlecht, nur hier das falsche Werkzeug.}
        Für grobe \emph{Flächen}-Aufgaben (z.\,B.\ Eisflächen-Aufnahmen von
        oben/unten) kann SAM durchaus geeignet sein.
  \item \textbf{Fine-Tuning von SAM unnötig.} Der flächenorientierte Decoder und die
        1024-Skalierung bleiben --- das Grundproblem wird dadurch nicht gelöst.
\end{enumerate}

\section{Geeignetere Modelle}
\label{sec:andere}
Die für dünne, gekrümmte Laserlinien passende Modellklasse sind
\emph{Segmentierer für kurvilineare Strukturen} --- genau die Familie, in der das
eigene U-Net liegt. Tabelle~\ref{tab:andere} ordnet ein.

\begin{table}[h]
  \centering
  \caption{Modelle/Ansätze und ihre Eignung für die 1-px-Laserlinie.}
  \label{tab:andere}
  \begin{tabular}{lll}
    \toprule
    \textbf{Modell / Ansatz} & \textbf{Typ} & \textbf{Eignung} \\
    \midrule
    SAM / SAM2 & Foundation, flächenorientiert & schlecht (hier getestet) \\
    HED, DexiNed & gelernte Kantendetektion & mittel (Kanten, nicht laser-spezifisch) \\
    Gefäß-/Riss-Segmentierung & U-Net-Varianten für dünne & \textbf{gut} (richtige Familie \\
    \;(DRIVE, DeepCrack) & Kurvenstrukturen & \;= unser U-Net) \\
    Steger-Liniendetektor & klassisch, sub-pixel & \textbf{sehr gut} (Vorverarbeitung) \\
    Frame-Diff / Canny & klassisch, voll aufgelöst & gut, schnell \\
    \textbf{U-Net from scratch (eigen)} & spezialisierte Segmentierung & gut, live-tauglich \\
    \bottomrule
  \end{tabular}
\end{table}

\noindent
Vortrainierte Encoder (z.\,B.\ ResNet über \texttt{segmentation\_models\_pytorch})
sind denkbar, bringen aber wenig: Die Domäne ist fern von ImageNet und die Aufgabe
ist \emph{low-level} (heller dünner Streifen). Das low-level-Vorwissen (Kanten,
Flecken) überträgt sich zwar, die objektspezifischen tiefen Merkmale sind irrelevant
--- bei mehr Rechenlast. Ein kompaktes U-Net from scratch genügt und ist schneller.

\section{Fazit}
SAM belegt \emph{quantitativ}, warum ein generisches Foundation-Modell an der
1-px-Laserlinie scheitert (feste 1024-Skalierung + flächen- statt kurvenorientierter
Decoder: 34\,\% Bildabdeckung statt $<1\,\%$). Das rechtfertigt sowohl die
klassischen, voll auflösenden Verfahren als auch ein spezialisiertes, von Grund auf
trainiertes U-Net zur semantischen Segmentierung der Laserlinie. Der Mehrwert von
SAM in dieser Arbeit ist damit \emph{methodisch}: eine ehrliche Baseline, die die
Modellwahl absichert.

\end{document}
